\chapter{System Design and Implementation}

This chapter presents the design and implementation of \textbf{Fiber}, the
end-to-end system that summarizes and fact-checks Vietnamese news. After a
requirements analysis (Section~3.1), Section~3.2 introduces the three-service
architecture; Sections~3.3--3.5 describe the summarization and
fact-checking modules; Section~3.6 details the three-axis evaluation
framework that is the methodological contribution of the thesis;
Section~3.7 covers the database schema; and Section~3.8 walks through the
user interface.

\section{Requirement Analysis}

\subsection{Functional Requirements}

\begin{enumerate}
  \item \textbf{Summarization.} Given a Vietnamese news URL, return an
        abstractive summary $S$ such that $|S| \leq 0.5 \cdot |D|$.
  \item \textbf{Streaming.} Stream the summary token-by-token over a
        Server-Sent-Events (SSE) channel for a sub-second time-to-first-token
        user experience.
  \item \textbf{Fact-checking.} On user request, decompose the summary into
        atomic claims, retrieve web evidence per claim, and render a trust
        score and source list.
  \item \textbf{Multi-model switching.} Allow the user to select one of
        fourteen supported LLMs at runtime; persist the selection across
        sessions.
  \item \textbf{Evaluation metrics (Axis~A).} After every summary, compute
        ROUGE-1/2/L, BLEU, BERTScore, compression rate, and write a row to
        the \texttt{evaluation\_metrics} table.
  \item \textbf{MoA fusion.} Provide a \textit{fusion} routing mode that
        runs three proposers in parallel and combines them through a
        GPT-4o aggregator.
  \item \textbf{Human-eval admin.} Allow the researcher to build $K$-way
        ranking tasks (model identities hidden) and to view aggregated
        Axis~C statistics.
\end{enumerate}

\subsection{Non-functional Requirements}

\begin{itemize}
  \item \textbf{Latency.} Time-to-first-token under one second for
        streaming summarization on a fibre connection; full summary in
        under five seconds for an article of $\leq 1000$ words.
  \item \textbf{Scalability.} Stateless backend with horizontal scaling
        via Vercel; BERTScore microservice deployable on Hugging Face
        Spaces.
  \item \textbf{API-key security.} All provider keys live on the server;
        the extension never sees an API key.
  \item \textbf{Cross-browser compatibility.} Manifest V3, validated on
        Chrome and Edge.
  \item \textbf{Cost tracking.} Every LLM call records its cost in USD
        to the metrics table for budget accounting.
\end{itemize}

\section{System Architecture}

Fiber is a three-service architecture (Figure~\ref{fig:architecture}).
The browser extension communicates with a Next.js backend over REST and
SSE; the backend talks to commercial LLM providers and the BERTScore
microservice; all persistent state lives in Supabase PostgreSQL.

\begin{figure}[H]
  \centering
  % TODO(figure): system architecture diagram --- see note at end of chapter.
  \fbox{\parbox{0.85\linewidth}{\centering Placeholder: three-service
  architecture diagram (Extension $\rightarrow$ Next.js backend
  $\rightarrow$ BERTScore microservice; Supabase as the shared
  persistence layer; LLM providers and Tavily as external endpoints).}}
  \caption{Fiber three-service architecture.}
  \label{fig:architecture}
\end{figure}

\subsection{Browser Extension (Plasmo)}

The extension is implemented in TypeScript with Plasmo and React. Three
runtime contexts are involved:

\begin{itemize}
  \item \textit{Content script.} Runs in the page world; extracts the
        article via Mozilla Readability, injects a Shadow-DOM-isolated
        sidebar, and forwards user actions to the background.
  \item \textit{Popup.} Compact UI for model switching, settings, and
        history.
  \item \textit{Background service worker.} Routes user actions to the
        backend, manages the SSE stream, and persists session state via
        \texttt{chrome.storage}.
\end{itemize}

\subsection{Backend API (Next.js 14)}

The backend uses the Next.js App Router. Each \texttt{/api/...} route
validates input with a Zod schema (\texttt{backend/domain/schemas.ts}),
delegates to a service in \texttt{backend/services/}, and returns a
strictly-typed JSON or SSE response. The Zod schemas are
single-source-of-truth contracts for both client and server, eliminating
a class of integration bugs.

Major routes:

\begin{itemize}
  \item \texttt{/api/summarize} --- LLM summarization, streaming or batch.
  \item \texttt{/api/fact-check} --- search-augmented verification.
  \item \texttt{/api/metrics} --- CRUD on evaluation rows.
  \item \texttt{/api/human-eval/*} --- Axis~C task and response endpoints.
  \item \texttt{/api/settings/*} --- routing, judge, and model configuration.
  \item \texttt{/api/logs/stream} --- SSE debug feed for development.
\end{itemize}

\subsection{BERTScore Microservice (FastAPI)}

BERTScore is computed in Python via the official \texttt{bert\_score}
library backed by \texttt{vinai/phobert-base}. The microservice exposes a
single \texttt{POST /calculate-score} endpoint and is deployed on
Hugging Face Spaces with a Docker container. Inputs are truncated to
$256$ tokens before scoring; the truncation is performed server-side so
that the contract from the Node backend is just ``send the two strings.''

\section{Text Summarization Module}

\subsection{Multi-Provider LLM Abstraction}

Fourteen models are supported across three providers
(\texttt{openai}, \texttt{@google/generative-ai},
\texttt{@anthropic-ai/sdk}). The \texttt{llm.service.ts} module exposes a
single \texttt{summarize} function that dispatches on provider and
applies provider-specific quirks:

\begin{itemize}
  \item Reasoning models (\texttt{o3-mini}, \texttt{o4-mini}) skip
        temperature, top-$p$, and frequency-penalty parameters.
  \item Anthropic temperature is clamped to $[0, 1]$; no seed support.
  \item Gemini uses \texttt{responseMimeType: "application/json"} for
        structured output.
\end{itemize}

Per-call cost is computed from the provider's published price list and
stored in the \texttt{cost\_usd} column of \texttt{evaluation\_metrics}.
Only one model can be ``active'' at a time, enforced by a partial unique
index in \texttt{model\_configurations}.

\subsection{Complexity-Based Routing}

For short Vietnamese inputs, an in-house Vietnamese model is cheaper and
more idiomatic than calling a frontier LLM. The routing layer estimates
input length as $\lceil |D|_{\text{chars}} / 4 \rceil$ tokens and dispatches
as follows:

\begin{itemize}
  \item $\leq 400$ tokens $\Rightarrow$ ViT5-large.
  \item $400$--$2000$ tokens $\Rightarrow$ GPT-4o-mini (default).
  \item $> 2000$ tokens $\Rightarrow$ GPT-4o-mini with chunked context.
\end{itemize}

Three \textit{routing modes} are exposed:

\begin{itemize}
  \item \texttt{auto} --- complexity-based dispatch as above.
  \item \texttt{forced} --- the user-selected model is always used.
  \item \texttt{fusion} --- the MoA pipeline (Section~3.4) is used.
\end{itemize}

\subsection{Streaming vs.\ Batch}

The streaming path uses Server-Sent Events. The backend opens an SSE
connection and forwards provider stream chunks as
\texttt{data:} frames; the extension renders tokens as they arrive.
Evaluation-metrics computation (ROUGE/BLEU/BERTScore) is dispatched
\textit{after} the response stream closes, in a fire-and-forget Promise
that writes to Supabase asynchronously; this keeps user-facing latency
independent of evaluation cost.

\section{MoA Output Fusion Pipeline}

The MoA pipeline is the main technical contribution of the thesis. It is
a faithful instantiation of Wang et al.\ \cite{Wang2024} with one
deliberate adaptation: a \textit{source-article residual connection}
into the aggregator prompt, which is necessary because the task is
grounded summarization rather than open-ended instruction-following.

\subsection{Two-Layer Design}

Layer~1 (proposers) is the set $\{\texttt{gpt-4o-mini},
\texttt{claude-haiku-4-5}, \texttt{gemini-2.5-flash}\}$. Each proposer
receives the same prompt --- ``summarize this Vietnamese news article in
five to seven sentences'' --- and produces a draft $d_i$ independently
(in parallel via \texttt{Promise.all}). Layer~2 is the aggregator
\texttt{gpt-4o}, which receives the source $D$ and the draft list
$\{d_1, d_2, d_3\}$ and produces the fused summary $S^{\star}$.

\begin{figure}[H]
  \centering
  % TODO(figure): MoA pipeline diagram --- see note at end of chapter.
  \fbox{\parbox{0.85\linewidth}{\centering Placeholder: MoA pipeline
  diagram --- three proposers in parallel feeding a single aggregator,
  with a source-article residual edge drawn from the input directly into
  the aggregator (Equation~1 of \cite{Wang2024}).}}
  \caption{MoA two-layer pipeline with source-article residual
  connection.}
  \label{fig:moa}
\end{figure}

\subsection{Proposer Layer}

The three proposers are commercial cloud LLMs from three different
families (OpenAI, Anthropic, Google) to maximise model diversity at low
unit cost (all three are the \textit{mini}/\textit{haiku}/\textit{flash}
tier of their families). Failure handling: if any one proposer times out
or returns an error, the remaining drafts are still fused; if two or
three proposers fail, the pipeline returns the surviving draft directly
and logs a warning.

\subsection{Aggregator Layer}

The aggregator prompt is the most important single artefact of the
implementation; it determines whether the aggregator merely stitches
drafts or actually synthesizes. The prompt, after the
2026-05-08 fix (commit \texttt{afac46e}), instructs the model to
\begin{itemize}
  \item evaluate the drafts critically rather than copy them verbatim;
  \item maintain a neutral Vietnamese journalistic register;
  \item produce a summary consistent with the source article;
  \item \textit{not} include the phrase ``cô đọng''
        (``condense / concise''), which empirically pushed the model
        toward truncated, draft-like outputs.
\end{itemize}

\subsection{Source-Article Residual Connection}

Equation~1 of \cite{Wang2024} concatenates the prompt with the
draft list. For grounded summarization, the prompt alone is insufficient
--- the aggregator needs the article to ground its synthesis. The
implementation injects the source article into the aggregator's system
context as an additional input, alongside the draft list. This is the
sole \textit{domain adaptation} of the paper's recipe.

The residual connection is empirically decisive. Without the article,
the aggregator falls back to stitching tokens from the drafts (compression
$\approx 22\%$ longer than average draft length, ROUGE-1 against the
source $\approx 0.34$); with the article, it produces editorial
synthesis ($\approx 17\%$ shorter, ROUGE-1 falls to $\approx 0.24$ on
strict prompts but holistic rubric scores rise). The full ablation is
reported in Section~4.2.

\subsection{Aggregator Prompt Ablation}

Three prompt variants were tested on 50 \textit{tienphong.vn} articles:

\begin{itemize}
  \item \textbf{Baseline} --- pre-fix prompt; no source in aggregator
        context.
  \item \textbf{v1 (strict source)} --- source injected, strict
        ``do not invent'' rules.
  \item \textbf{v2 (soft source)} --- source injected, softer
        ``use as reference'' wording.
\end{itemize}

The headline finding (reported in detail in Section~4.2) is that the
mere \textit{presence} of the source article in the aggregator's context
flips behavior; the precise wording of the rules around it does not.
The shipped prompt is v2 with the post-\texttt{afac46e} removal of the
``cô đọng'' phrase.

\section{Fact-Checking Module}

The fact-check module is a single LLM-orchestrated RAG instance.

\begin{enumerate}
  \item The summary (or an arbitrary user-pasted snippet) is sent to
        Tavily as a search query.
  \item Tavily returns six results, each a
        \texttt{(title, url, snippet)} triple.
  \item The snippets, together with the original claim text, are passed
        to GPT-4o-mini, which returns a JSON object with
        \texttt{trust\_score} ($0$--$100$), \texttt{verified} (boolean),
        and \texttt{sources} (list of URLs).
  \item The extension renders this in a modal beneath the summary.
\end{enumerate}

The choice of GPT-4o-mini for the verdict step is deliberate: it is
ten times cheaper than GPT-4o and shows no measurable drop in
verdict quality on a $50$-claim spot-check.

\section{Three-Axis Evaluation Framework}

This section describes the methodological contribution: a single
evaluation framework that combines content-retention metrics
(Axis~A), LLM-judge quality (Axis~B), and blind human ranking
(Axis~C). The choice of three axes follows from the analysis in
Section~2.2.5: any one axis alone is structurally biased; together
they triangulate.

\subsection{Axis A --- Content Retention}

Axis~A is computed in
\texttt{backend/services/evaluation.service.ts} for every summary.
ROUGE-1/2/L and BLEU use the \texttt{rouge} and \texttt{sacrebleu}
JavaScript ports. BERTScore is delegated to the Python microservice
(\texttt{bert.service.ts}). All four metrics are computed against the
\textit{source article}, not a human reference summary; this is
documented in the unified report as a methodology caveat. The
compression rate is $|S|/|D|$ in characters.

\subsection{Axis B --- Quality and Preference}

Axis~B is the headline axis of the thesis. The
\texttt{llm-judge.service.ts} module exposes four functions:

\begin{itemize}
  \item \texttt{judgeRubric} --- FLASK-derived 5-dimension rubric on a
        1--5 scale.
  \item \texttt{judgeAbsolute} --- MT-Bench-style 1--10 holistic score.
  \item \texttt{judgePairwise} --- A vs.\ B with position randomization.
  \item \texttt{scoreFactuality} (in \texttt{factuality.service.ts})
        --- claim split + entailment classification.
\end{itemize}

Two higher-level routines operate on MoA fusion outputs:

\begin{itemize}
  \item \texttt{runFusionPairwiseJudge} --- fused vs.\ best-draft
        ($\texttt{comparison\_type} = \texttt{vs\_best\_draft}$).
  \item \texttt{runFusionVsAllDraftsJudge} --- fused vs.\ each
        individual proposer
        ($\texttt{comparison\_type} = \texttt{vs\_individual\_draft}$).
  \item A third comparison, ``fused vs.\ \texttt{gpt-4o} running
        alone'' ($\texttt{vs\_single\_aggregator}$), is generated by an
        offline script described in Section~4.3.2; this is the
        thesis-decisive comparison because both candidates are
        \texttt{gpt-4o} output and the same-family judge bias cancels.
\end{itemize}

\subsection{Axis C --- Human Validation}

Axis~C uses the same extension's web UI to administer blind
ranking tasks.

\begin{itemize}
  \item The researcher creates a task at \texttt{/evaluate/admin}.
        Each task references one article and three candidate summaries
        (fused, GPT-4o-alone, GPT-4o-mini draft) with model identities
        hidden. The candidates are presented as
        \textit{Bản A, Bản B, Bản C}.
  \item A rater opens \texttt{/evaluate?task=<uuid>}, types a stable
        identifier (the ``rater code''), drag-ranks the three
        candidates, and writes a one-sentence Vietnamese rationale per
        candidate.
  \item Responses are written to \texttt{human\_eval\_responses}; a
        partial unique index on \texttt{(task\_id, rater\_id)}
        prevents double-submission.
  \item \texttt{/api/human-eval/report} aggregates per-approach
        average rank, per-approach win rate, and Fleiss'~$\kappa$.
\end{itemize}

\subsection{Statistical Analysis}

All statistical primitives live in
\texttt{backend/output-fusion/scripts/stats.ts}:

\begin{itemize}
  \item \texttt{signTestPValue(wins, losses, ties)} --- exact two-sided
        binomial.
  \item \texttt{lengthBucketedWinRate(pairs)} --- Dubois-2024 simplified
        bucket method with $\text{MIN\_BUCKET\_N} = 5$.
  \item \texttt{fleissKappa(matrix)} and
        \texttt{fleissKappaFromRankings(rankings)} --- inter-rater
        agreement.
  \item \texttt{pairedMetricStats(a, b)} --- mean, stdev, and paired
        differences for Axis-A side-by-side reporting.
\end{itemize}

\subsection{P0-8: Fused vs.\ GPT-4o-Alone (Thesis-Decisive)}

The decisive comparison removes a confound that plagues fused
vs.\ \textit{best-proposer} comparisons: the best proposer is a
weaker model than the aggregator, so any fused-vs-best win could be
attributed to ``GPT-4o is just better than Haiku''. By comparing
fused (which uses GPT-4o as aggregator) against GPT-4o running alone
(zero-shot, same single article), both candidates are GPT-4o output.
A win is therefore attributable to the act of
\textit{synthesis} --- giving the aggregator multiple drafts as
context --- not to model-capability gap. The implementation is
\texttt{scripts/run-single-baseline.ts} +
\texttt{scripts/compare-fused-vs-single.ts}; results in Section~4.3.2.

\section{Database Design}

\subsection{Schema Overview}

The Supabase PostgreSQL database holds twenty-five migrations. The
schema is organised into four functional groups:

\begin{itemize}
  \item \textbf{Per-summary metrics:} \texttt{evaluation\_metrics},
        \texttt{routing\_decisions}.
  \item \textbf{MoA-specific:} \texttt{moa\_fusion\_results},
        \texttt{moa\_draft\_results},
        \texttt{llm\_judge\_pairwise}.
  \item \textbf{Human evaluation:} \texttt{human\_eval\_tasks},
        \texttt{human\_eval\_responses}.
  \item \textbf{Configuration:} \texttt{app\_settings},
        \texttt{model\_configurations},
        \texttt{dashboard\_actions}.
\end{itemize}

\subsection{Key Tables}

\textbf{\texttt{evaluation\_metrics}.} One row per summary. Holds Axis~A
columns (\texttt{rouge1\_f}, \texttt{rouge2\_f}, \texttt{rougel\_f},
\texttt{bleu}, \texttt{bertscore\_f1}, \texttt{compression\_pct},
\texttt{latency\_ms}) and Axis~B columns
(\texttt{judge\_rubric} JSONB, \texttt{judge\_absolute},
\texttt{judge\_cost\_usd}, \texttt{factuality\_entailed\_ratio},
\texttt{factuality\_hallucinations} JSONB).

\textbf{\texttt{llm\_judge\_pairwise}.} One row per pairwise verdict.
The \texttt{comparison\_type} column is a discriminator with three
values: \texttt{vs\_best\_draft} (Axis B.2),
\texttt{vs\_individual\_draft} (Axis B.2b, one row per (fusion,
proposer) pair), and \texttt{vs\_single\_aggregator}
(Axis B.2c, P0-8).

\textbf{\texttt{moa\_fusion\_results} / \texttt{moa\_draft\_results}.}
One row per fusion run; one row per proposer draft. The
\texttt{pipeline\_mode} column is currently always
\texttt{moa\_synthesis}; the historical \texttt{llm\_ranker} value
remains for rows written before 2026-05-09.

\textbf{\texttt{human\_eval\_tasks} / \texttt{human\_eval\_responses}.}
Tasks carry a hidden \texttt{candidates} JSONB array with model
labels; rater responses carry a ranking (list of label letters) and a
per-candidate Vietnamese \texttt{rationale}. The unique index
\texttt{(task\_id, rater\_id)} prevents duplicate submissions.

\section{User Interface}

\subsection{Extension Sidebar}

The extension sidebar is the primary user surface
(Figure~\ref{fig:sidebar}). Tapping the toolbar icon injects a
right-edge panel into the current page; the panel shows a
streaming-summary area, a fact-check button (which opens the modal
described in Section~3.5), a model selector, and a kebab menu with
``copy summary'', ``open metrics dashboard'', and ``settings''.

\begin{figure}[H]
  \centering
  % TODO(figure): screenshot of the extension sidebar.
  \fbox{\parbox{0.85\linewidth}{\centering Placeholder: screenshot of
  the Fiber sidebar on a \textit{tienphong.vn} article, showing the
  streaming summary, fact-check button, and model selector.}}
  \caption{The Fiber extension sidebar.}
  \label{fig:sidebar}
\end{figure}

\subsection{Metrics Page (\texttt{/metrics})}

The metrics page renders the \texttt{evaluation\_metrics} table with
two view modes (\textit{compact}, \textit{full}), filters for routing
mode and model, and three inline widgets: a five-axis rubric radar,
a fused-vs-best-draft pairwise badge, and a factuality entailment
ring.

\subsection{Evaluation Admin (\texttt{/evaluate/admin})}

The admin page has two tabs:

\begin{itemize}
  \item \textit{Create.} Pick a source article URL and three candidate
        summaries from the metrics list; the page builds a hidden-label
        task and prints a shareable
        \texttt{/evaluate?task=<uuid>} link.
  \item \textit{Review.} Aggregates responses per task and across the
        cohort: per-approach average rank, per-approach win rate, and
        Fleiss'~$\kappa$ with the Landis--Koch band label.
\end{itemize}

\subsection{Rater UI (\texttt{/evaluate})}

The rater UI is Vietnamese-only and presents the three candidates as
draggable cards labelled \textit{Bản A}, \textit{Bản B},
\textit{Bản C} above the source article. A drag-and-drop ordering
mechanism (with $\blacktriangle\,\blacktriangledown$ accessibility
buttons) records the ranking; a textarea beneath each card records the
rationale. The header is suppressed so that hidden labels stay hidden.

\subsection{Settings (\texttt{/settings})}

The settings page exposes three configuration cards:

\begin{itemize}
  \item \textbf{Routing.} Mode (\texttt{auto} / \texttt{forced} /
        \texttt{fusion}) and forced-model selector.
  \item \textbf{Judge.} Mode (\texttt{metrics\_only} /
        \texttt{judge\_only} / \texttt{both}), default judge model,
        default style (rubric / absolute), and a factuality toggle.
  \item \textbf{Provider keys (read-only).} Reports which provider
        API keys are configured server-side; the keys themselves
        never reach the client.
\end{itemize}

\section{Summary}

This chapter described the design and implementation of Fiber along
three threads: the production system (extension + Next.js backend +
BERTScore service), the MoA fusion pipeline (the technical
contribution), and the three-axis evaluation framework (the
methodological contribution). The pipeline is faithful to
\cite{Wang2024} apart from the source-article residual connection,
which is necessary for grounded summarization. The three-axis
framework reads measurements through three structurally-different
lenses, so that any conclusion supported by all three is unlikely to
be an artefact of any single metric.

\vspace{0.8em}
\noindent\textit{Figures needed for this chapter.} (1)~Figure~3.1, the
three-service system architecture; (2)~Figure~3.2, the MoA pipeline
diagram with the source-article residual connection;
(3)~Figure~3.3, a screenshot of the extension sidebar. All three are
flagged with placeholders in this chapter.
